%==============================================================================
%  CHAPTER 5 — STUDY AND IMPLEMENTATION OF SPRINT 3
%==============================================================================
\chapter{Study and Implementation of Sprint 3}
\label{chap:sprint3}

\section*{Introduction}
The third sprint, \emph{Model Registry and Deployment}, turns tracked runs into
deployable assets. It adds the model registry with lifecycle stages, one-click
deployment of a model as a live service, and an in-interface prediction tester.

% ---------------------------------------------------------------------------
\section{Sprint Backlog}
\begin{table}[H]\centering
\renewcommand{\arraystretch}{1.25}
\caption{Sprint~3 backlog}\label{tab:backlog-s3}
\begin{tabularx}{\textwidth}{|c|X|c|}
\hline
\textbf{\#} & \textbf{Task} & \textbf{Pts}\\
\hline
8  & Register model versions from successful runs & 2\\ \hline
8  & Build the model registry interface & 3\\ \hline
9  & Promote models across Staging/Production/Archived & 3\\ \hline
10 & Deploy a model as an InferenceService on the cluster & 5\\ \hline
10 & Track deployment status and endpoint & 3\\ \hline
11 & Build the in-UI prediction tester & 5\\ \hline
\end{tabularx}
\end{table}

% ---------------------------------------------------------------------------
\section{Analysis and Functional Specification}

\subsection{Use Case Diagram of Sprint 3}
This sprint covers \emph{Register a model}, \emph{Promote a model},
\emph{Deploy a model}, and \emph{Test predictions}. Figure~\ref{fig:uc-s3} shows
the use-case diagram.
\begin{figure}[H]\centering
  \includegraphics[width=0.9\textwidth,height=0.72\textheight,keepaspectratio]{uc_sprint3}
  \caption{Use-case diagram of Sprint~3}
  \label{fig:uc-s3}
\end{figure}

\subsection{Use Case Refinement of Sprint 3}
Tables~\ref{tab:uc-deploy} and~\ref{tab:uc-promote} refine the \emph{Deploy a
model} and \emph{Promote a model} use cases.
\begin{table}[H]\centering
\renewcommand{\arraystretch}{1.3}
\caption{Textual description of the ``Deploy a model'' use case}\label{tab:uc-deploy}
\begin{tabularx}{\textwidth}{|>{\bfseries}p{3.2cm}|X|}
\hline
Use case & Deploy a model\\ \hline
Actor & Platform user\\ \hline
Precondition & A model version is registered in the registry.\\ \hline
Post-condition & A live prediction service is running and reachable on the cluster.\\ \hline
Nominal scenario & 1.~The user clicks ``Deploy'' on a registered model.\newline 2.~The backend creates an \texttt{InferenceService} resource.\newline 3.~The serving operator reconciles it into a Deployment, Service, and Ingress.\newline 4.~The model server loads the artifact from object storage.\newline 5.~The deployment status and endpoint are reported back to the user.\\ \hline
Exceptions & The service fails to become ready: the deployment is marked failed and the user is notified.\\ \hline
\end{tabularx}
\end{table}

\begin{table}[H]\centering
\renewcommand{\arraystretch}{1.3}
\caption{Textual description of the ``Promote a model'' use case}\label{tab:uc-promote}
\begin{tabularx}{\textwidth}{|>{\bfseries}p{3.2cm}|X|}
\hline
Use case & Promote a model\\ \hline
Actor & Platform user\\ \hline
Precondition & A model version is registered in the registry.\\ \hline
Post-condition & The model version carries its new lifecycle stage in both the registry and the platform.\\ \hline
Nominal scenario & 1.~The user selects a target stage (Staging, Production, or Archived).\newline 2.~The system transitions the version in the model registry.\newline 3.~The system mirrors the new stage in its own database.\newline 4.~The stage badge updates in the interface.\\ \hline
Exceptions & The registry is unreachable: no partial state is written and the user is informed.\\ \hline
\end{tabularx}
\end{table}

% ---------------------------------------------------------------------------
\section{Design}
This sprint adds two entities, described in Table~\ref{tab:entities-s3}:
\texttt{MLModel}, the platform's mirror of a registered model version, and
\texttt{Deployment}, the record of a model running as a live service.

\begin{table}[H]\centering
\renewcommand{\arraystretch}{1.25}
\caption{Entities introduced in Sprint~3}\label{tab:entities-s3}
\begin{tabularx}{\textwidth}{|l|l|X|}
\hline
\textbf{Entity} & \textbf{Attribute} & \textbf{Description}\\
\hline
\multirow{4}{*}{\texttt{MLModel}}
 & id & Unique identifier (primary key)\\ \cline{2-3}
 & project\_id & Reference to the owning \texttt{Project}\\ \cline{2-3}
 & name, version & Registered model name and version number\\ \cline{2-3}
 & stage & Staging, Production, or Archived\\
\hline
\multirow{4}{*}{\texttt{Deployment}}
 & id & Unique identifier (primary key)\\ \cline{2-3}
 & model\_id & Reference to the deployed \texttt{MLModel}\\ \cline{2-3}
 & inference\_service\_name & Name of the cluster resource\\ \cline{2-3}
 & status, endpoint & Readiness state and prediction \gls{url}\\
\hline
\end{tabularx}
\end{table}

\subsection{Sequence Diagrams}
Figure~\ref{fig:seq-promote} details the promotion flow, and
Figure~\ref{fig:seq-deploy} the deploy-and-predict flow.
\begin{figure}[H]\centering
  \includegraphics[width=0.9\textwidth,height=0.62\textheight,keepaspectratio]{seq_promote}
  \caption{Sequence diagram of promoting a model version}
  \label{fig:seq-promote}
\end{figure}
\begin{figure}[H]\centering
  \includegraphics[width=0.9\textwidth,height=0.62\textheight,keepaspectratio]{seq_deploy}
  \caption{Sequence diagram of deploying a model and requesting a prediction}
  \label{fig:seq-deploy}
\end{figure}

% ---------------------------------------------------------------------------
\section{Implementation}

\subsection{Model Registry and Promotion}
Successful runs can be registered as named, versioned models. Each version carries
a lifecycle stage, Staging, Production, or Archived, which the user changes
from the interface. The tracking server remains the source of truth for the
registry; the platform mirrors and controls it through its own \gls{ui}.

\subsection{Deployment as an InferenceService}
Deploying a model reduces to creating a single declarative resource on the cluster.
The serving operator reconciles this resource into a running deployment, a service,
and an ingress, backed by a framework-specific model server that loads the model
artifact directly from the object store.

\begin{lstlisting}[language=Python,caption={Deployment is expressed declaratively; the operator does the rest.},label={lst:isvc}]
inference_service = {
    "apiVersion": "serving.kserve.io/v1beta1",
    "kind": "InferenceService",
    "metadata": {"name": name,
                 "annotations": {"serving.kserve.io/deploymentMode": "RawDeployment"}},
    "spec": {"predictor": {"sklearn": {"storageUri": model_artifact_uri}}},
}
\end{lstlisting}

\subsection{Deployment Status Tracking}
An \texttt{InferenceService} does not become ready instantly: the operator must
create the deployment, the model server must start, and the artifact must be
downloaded from object storage. The backend therefore polls the resource's status
conditions and updates the \texttt{Deployment} row as the service progresses,
so the interface always reflects the real state of the cluster.

\begin{lstlisting}[language=Python,caption={Reading readiness from the InferenceService status conditions.},label={lst:isvc-status}]
def isvc_status(isvc: dict) -> str:
    for cond in isvc.get("status", {}).get("conditions", []):
        if cond["type"] == "Ready":
            return "READY" if cond["status"] == "True" else "PENDING"
    return "PENDING"
\end{lstlisting}

\subsection{Prediction}
The backend forwards prediction requests to the model server through the cluster's
\gls{api}, which avoids needing a publicly-routable address for every model. The
interface provides a tester where the user submits sample inputs and sees the
returned predictions.

\subsection{User Interface}
The interface shows the registry with stage badges, a deployment panel reporting
status and endpoint, and a prediction tester
(Figures~\ref{fig:screen_registry} to~\ref{fig:screen_predict}).
\screenshot{screen_registry}{Model registry with lifecycle stage badges}
\screenshot{screen_deploy}{Deployment panel reporting status and endpoint}
\screenshot{screen_predict}{In-interface prediction tester}

% ---------------------------------------------------------------------------
\section{Testing and Validation}
Table~\ref{tab:tests-s3} summarizes the main test scenarios of Sprint~3.

\begin{table}[H]\centering
\renewcommand{\arraystretch}{1.25}
\caption{Sprint~3 test scenarios}\label{tab:tests-s3}
\begin{tabularx}{\textwidth}{|X|X|c|}
\hline
\textbf{Scenario} & \textbf{Expected result} & \textbf{Status}\\
\hline
Register a model from a successful run & Version appears in the registry with stage None & OK\\ \hline
Promote a version to Production & Stage updated in registry and interface & OK\\ \hline
Deploy a registered model & InferenceService created; status progresses to READY & OK\\ \hline
Deploy with an invalid artifact path & Deployment marked failed; user notified & OK\\ \hline
Send a prediction from the tester & Prediction returned from the model server & OK\\ \hline
Delete a deployment & Cluster resources removed; record cleaned up & OK\\ \hline
\end{tabularx}
\end{table}

% ---------------------------------------------------------------------------
\section{Sprint Review}
At the end of Sprint~3 a user can take a tracked run all the way to a live,
queryable model, the platform now \emph{registers, deploys, and serves}.
\begin{figure}[H]\centering
\begin{tikzpicture}
\begin{axis}[
    width=0.8\textwidth, height=6cm,
    xlabel={Day}, ylabel={Remaining points},
    xmin=0, xmax=11, ymin=0, ymax=22,
    legend pos=north east, grid=both,
]
\addplot[dashed,gray] coordinates {(0,21)(10,0)};
\addplot[thick,primary,mark=*] coordinates {
  (0,21)(1,21)(2,19)(3,17)(4,14)(5,12)(6,9)(7,7)(8,4)(9,2)(10,0)};
\legend{Ideal, Actual}
\end{axis}
\end{tikzpicture}
\caption{Sprint~3 burndown chart}\label{fig:burndown-s3}
\end{figure}

\noindent\textbf{Retrospective.}
What went well: expressing deployment as a single declarative resource paid off,
once the operator was configured, deploying any model became trivial and uniform.
To improve: configuring the serving operator to match the cluster's ingress setup
(RawDeployment mode instead of the default serverless mode) required careful
investigation and took several days of the sprint; reading the operator's
documentation end-to-end before starting would have surfaced the constraint
earlier.

\section*{Conclusion}
Sprint~3 closed the loop from a tracked run to a live, queryable model. The final
sprint makes the platform observable, opens it to external clients, and adds
administration.
